The measurement is performed using $\Z$ events, as the muons from $\JPsi$ decays need not be
compatible with the primary vertex if the $\JPsi$ is produced in the decay of a B hadron.
For this measurement, the probe is a muon passing the POG Loose identification criteria,
and it is considered a passing probe if satisfies the SIP3D, $d_{xy}$, $d_{z}$ cuts of this analysis.

For Run2, very good agreement between data and simulation is observed in the barrel (Fig.~\ref{fig:MuonIDEff_2}, left)
while some inefficiency is visible in the endcaps, especially at large values of $|\eta|$.
The data to simulation scale factor is found to be flat as function of \pt, so, similarly to what done
for the identification part, we apply a correction only as function of $\eta$.


%%%%%%%%%%
\begin{figure}[!h]
  \begin{center}
    \subfigure[]{\includegraphics[width=0.32\textwidth]{Figures/chapter3_Figure/mu_SIP4_barrel.pdf}}
    \subfigure[]{\includegraphics[width=0.32\textwidth]{Figures/chapter3_Figure/mu_SIP4_endcap.pdf}}
    \subfigure[]{\includegraphics[width=0.32\textwidth]{Figures/chapter3_Figure/mu_SIP4_pt20.pdf}}
    \caption{Efficiency of the muon impact parameter requirements, measured with the tag\&probe method on \Z events, as function of \pt in the barrel (left) and 
    endcaps (center), and as function of $\eta$ for $\pt > 20\GeV$ (right). In the upper panel, the larger error bars include also the systematical
    uncertainties, while the smaller ones are purely statistical. In the lower panel showing the ratio of the two efficiencies, the black error bars are 
    for the statistical uncertainty, the orange rectangles for the systematical uncertainty and the violet rectangles include both uncertainties.}
    \label{fig:MuonIDEff_2}
\end{center}
\end{figure}

%%%%%%%%%
\par
For Run1, the efficiency of the $|SIP3D| < 4$ criteria is found to be above
$99.5\%$ in the barrel, and decreasing slightly down to about $99\%$
in the forward part of the detector (Fig.~\ref{fig:mu-tnp-sip}). In
the latter region, the efficiency in data about $0.1\%$ lower than in
simulations. A worse $|SIP3D| < 4$ efficiency due to pixel misalignment
in the promptly reconstructed data is now recovered in the
re-reconstructed data as can be seen in Fig.~\ref{fig:mu-tnp-sip} (right bottom).
Morover, the efficiency does not show a significant
dependance to pile-up as it is flat as a function of the number of
vertices as shown in Fig.~\ref{fig:mu-tnp-sip} (left bottom).

\begin{figure}[htb]
\begin{center}
\includegraphics[width=0.45\linewidth]{Figures/chapter3_Figure/HZZ_2011_MuonSIP_eta.pdf}
\includegraphics[width=0.45\linewidth]{Figures/chapter3_Figure/HZZ_2012_MuonSIP_eta.pdf} \\
\includegraphics[width=0.45\linewidth]{Figures/chapter3_Figure/HZZ_2012_MuonSIP_vtx.pdf}
\includegraphics[width=0.45\linewidth]{Figures/chapter3_Figure/HZZ_2012_Compare_MuonSIP_eta.pdf}
\caption{Efficiency for the requirement on the 3D impact parameter
  significance $|SIP3D| < 4$. Top row: efficiency as as function of
  the muon pseudorapidity, for 2011 data (left), and 2012 data
  (right), for muons with $\pt > 20 \GeV$; bottom row: efficiency in
  2012 data as function of the number of reconstructed vertices (left) and
  the comparison between the 2012 prompt-reco and the re-reco (right). Muons satisfying the Particle Flow identification
  requirements are used.}
\label{fig:mu-tnp-sip}
\end{center}
\end{figure}
%%%%%%%
